
\section{Implementation}
\label{sec:implementation}

\subsection{Components Overview}

Our prototype extends Linux 6.15 with a minimal kernel core ($\sim$1 KLOC) for hook registration and verifier extension, alongside a lightweight integration into NVIDIA's open GPU kernel modules ($\sim$100 LOC) for memory and scheduling hooks (§4.1–§4.2). The userspace loader/JIT ($\sim$10 KLOC) extends bpftime's infrastructure, while our LLVM-based backend ($\sim$1 KLOC) handles the translation of the restricted \sys eBPF subset into device code (PTX/SPIR‑V).

\subsection{Host-side Integration}
\label{sec:impl-host}

\subsubsection{Driver hooks for memory and scheduling}

\sys builds on the NVIDIA Open GPU Kernel Modules (Turing, Ampere, Hopper, and Blackwell or newer). We use the UVM module for memory operations, as it exposes critical memory decision points such as page fault handlers and prefetch logic with a mechanism to offload and manage memory on the CPU. This allows us to customize the kernel module behavior instead of letting users bypass all kernel mechanisms. Unlike prior approaches that manage UVM as a monolithic region, \sys enforces policy at the granularity of individual allocations or pages, implementing the region-level Add/Access/Remove and prefetch events described in Section~\ref{sec:design-mem}. Each callback constructs a compact context (e.g., region range, fault address, block size) and invokes an attached eBPF handler. The handler 
returns compact decisions (default or bypass), allowing for custom logic or the use of the kernel's default logic.

For scheduling, we extend the kernel lifecycle code to invoke \texttt{struct\_ops} callbacks at submission, admission, preemption, and completion points, allowing policies to inspect kernel configurations (e.g., grid size), reorder ready queues, or update per-tenant statistics dynamically.

On the host, \sys reuses the standard kernel eBPF verifier. We implement a kernel module that extends the verification context to support our GPU-specific attach points and special kfuncs. This module leverages the kernel's annotation mechanisms (BTF and kfunc
declarations) to enforce interface correctness and type safety, similar to the approach
in the bpftime paper~\cite{bpftime}. Where policies must trigger specific actions in the driver (such as committing a victim selection or marking a kernel for preemption), we expose these interfaces as safe kfuncs that can be called from eBPF under verifier control.

\subsubsection{\sysdevice user-space interface}

\sysdevice provides eBPF-style uprobe and uretprobe hooks for GPU-related
user-level functions, enabling bcc-style observability tools and higher-level
policy logic without modifying applications. Developers attach eBPF programs to
function entry and exit to collect metrics such as per-call latency, counts, or
arguments, using standard eBPF maps.

These hooks function as standard user-space probes on the host, but can be transparently offloaded to the device by the GPU runtime (see §\ref{sec:impl-gpu}), enabling high-frequency instrumentation without context switches.

\subsection{GPU Runtime}
\label{sec:impl-gpu}

\paragraph{Code generation and trampolines}
When a \sys program is verified and loaded into the kernel, our loader intercepts it, extracts the bytecode and BTF metadata, and hands them to the user-space backend to generate device code (PTX/SPIR-V). All execution remains in the kernel (via the normal eBPF JIT) or on the GPU via our backend and runtime. 

To enable dynamic instrumentation without recompilation, \sys performs binary rewriting at kernel launch, injecting small trampolines at strategic locations such as entry points, sampled memory operations, and phase boundaries.

A small set of runtime helpers with verifier-approved prototypes is exposed and pre-compiled to GPU device code. During translation, we perform static analysis to minimize context structures, materializing only used fields, and wrap handlers in prologue/epilogues that manage register preservation for the designated execution lane.
By driving device code generation directly from verified eBPF bytecode, \sys automatically extends the Linux kernel's memory and type safety guarantees to the GPU.


\paragraph{Scheduling and Memory Operations}

For scheduling policies, we implement a thread-block scheduler that models each
kernel as a collection of work units (e.g., tiles or logical block indices)
consumed by persistent worker blocks. On Blackwell NVIDIA GPUs, we map this to cluster
launch control\cite{}, which treats a grid of blocks as cancellable work units and enables true
hardware work-stealing style scheduling. Each worker processes its current unit, invokes a
device-side \sys handler (running in one lane) to select or approve the next work unit,
and uses the cluster API to claim or release indices. Progress and behavior probes are
implemented with \sys trampolines at kernel entry and annotated phase boundaries.
All device-side scheduling logic is subject to the same bounded-execution constraints as other \sys code. For memory policies, \sys attaches to specific instructions, tracepoints, or function boundaries to issue hints to the host via helper functions. Crucially, unlike eGPU prototypes that rely on heavy synchronization (fences and atomics), \sys helpers and attach points operate without explicit fences or atomics to minimize overhead.

\paragraph{Cross-domain State and Synchronization}

\sys uses eBPF maps as the unified state abstraction across host and device. We
allocate maps and coordination buffers in memory regions accessible to CPU and GPU,
avoiding explicit marshalling. On our NVIDIA prototypes we use managed shared
memory so that map contents are addressable in both contexts. Maps store policy
state such as per-region hotness, per-tenant scheduling priorities, per-block
progress, and counters for observability tools. Unlike eGPU, which placed the maps entirely on the CPU, causing prohibitive overhead, we allow maps to be placed inside GPU registers, GPU main memory, or CPU.
